\chapter{Limites du travail}

Ce projet remplit les objectifs pédagogiques et expérimentaux du module, mais plusieurs limites doivent être explicitées.

\begin{itemize}[leftmargin=1.2cm]
    \item Les budgets d’entraînement ont été volontairement limités pour préserver l’exécutabilité sur ordinateur portable.
    \item La partie tabulaire ne compare pas le MLP à des modèles non neuronaux souvent très compétitifs sur ce type de données.
    \item La partie vision reste sur Fashion-MNIST, dataset utile pédagogiquement mais modeste au regard de problèmes visuels modernes.
    \item La partie séquentielle repose sur un sous-corpus filtré, ce qui simplifie la tâche mais limite la couverture linguistique.
    \item Le système Seq2Seq n’intègre ni attention explicite, ni Transformer, ni embeddings pré-entraînés.
    \item Les résultats séquentiels doivent donc être interprétés comme une démonstration robuste de pipeline et de principes, non comme un système de traduction à l’état de l’art.
\end{itemize}

Ces limites n’invalident pas le projet ; elles en fixent simplement le périmètre scientifique et computationnel.

